% GENERATED FILE. Edit canonical lessons, then run npm run generate:books.
% canonical-source-hash: fnv1a64:81b41ea823a56122
% canonical-lessons: BN-C302-kahini, BN-C302-shonano, BN-C302-bornona-kora, BN-C302-ghotano, BN-C302-kromosh

\chapter{Telling It Well}
\label{ch:bn-telling-it-well}
\hlchaptermodality{302}

\begin{chapteropening}
\textbf{By the end of this chapter:} \emph{I can tell a story and say how it slowly unfolded.}

Complete the last lesson of chapter 302: I can tell a story and say how it slowly unfolded.
\end{chapteropening}

\section[\texorpdfstring{\bn{কাহিনি} (kāhini)}{কাহিনি (kāhini)}]{\bn{কাহিনি} (kāhini) — a tale; a narrative}

\label{lesson:BN-C302-kahini}



\begin{quote}
Before the new one: say the Bengali for detailed, then the Bengali for to verify.
\end{quote}

\subsection*{You'll want to know: \bn{কাহিনি}}
\textbf{\bn{কাহিনি}} — \emph{kāhini} — \textquotedblleft{}a tale; a narrative\textquotedblright{}.

\begin{quote}
\bn{দাদু} \bn{রাতে} \bn{আমাদের} \bn{একটা} \bn{পুরনো} \bn{কাহিনি} \bn{বললেন।}
\end{quote}

\begin{quote}
\emph{dādu rāte āmāder ekṭā purôno kāhini bôllen.}
\end{quote}

\begin{quote}
At night my grandfather told us an old tale.
\end{quote}

\begin{cousinweb}
\textbf{\bn{কাহিনি}} (\emph{kāhini}) comes down from Sanskrit \emph{kathana}, 'narration' (CDIAL 2702, which lists Bengali \emph{kahinī}, 'a story'); it is a cousin of \textbf{\bn{কথা}}, talk.
\end{cousinweb}

\begin{grammarlens}[title={What you've built}]
One more word for telling what happened.
\end{grammarlens}

\subsection*{Your turn}
\begin{itemize}
\raggedright
  \item \emph{Say it:} \emph{kāhini}
  \item \emph{Say it:} \emph{kāhini}, once more
  \item \emph{Say it:} \emph{dādu rāte āmāder ekṭā purôno kāhini bôllen.} — At night my grandfather told us an old tale.
  \item \emph{Recall:} say the Bengali for detailed, then the Bengali for to verify, then say \emph{kāhini} again
\end{itemize}

\begin{tcolorbox}[breakable,colback=teal!4,colframe=teal!35!black,title={Before you move on}]
What does \bn{কাহিনি} mean? (\textquotedblleft{}A tale; a narrative\textquotedblright{}.) Say it once more.
\end{tcolorbox}
\section[\texorpdfstring{\bn{শোনানো} (shonāno)}{শোনানো (shonāno)}]{\bn{শোনানো} (shonāno) — to tell a story to; to recite}

\label{lesson:BN-C302-shonano}



\begin{quote}
Before the new one: say the Bengali for to verify, then the Bengali for a tale.
\end{quote}

\subsection*{You'll want to know: \bn{শোনানো}}
\textbf{\bn{শোনানো}} — \emph{shonāno} — \textquotedblleft{}to tell a story to; to recite\textquotedblright{}.

\begin{quote}
\bn{দিদিমা} \bn{আমাদের} \bn{রাজা} \bn{আর} \bn{রানির} \bn{গল্প} \bn{শোনালেন।}
\end{quote}

\begin{quote}
\emph{didimā āmāder rājā ār rānir gôlpo shonālen.}
\end{quote}

\begin{quote}
Grandmother told us the story of a king and a queen.
\end{quote}

Past (he/she …): \bn{তিনি} \bn{একটা} \bn{গান} \bn{শোনালেন}, \emph{tini ekṭā gān shonālen}.

\begin{cousinweb}
\textbf{\bn{শোনানো}} (\emph{shonāno}) is \textbf{\bn{শোনা}}, to hear — CDIAL 12598, Sanskrit \emph{śṛṇōti}, 'hears' — with the -ānō ending that makes someone else hear: to tell a story to someone.
\end{cousinweb}

\begin{grammarlens}[title={What you've built}]
One more word for telling what happened.
\end{grammarlens}

\subsection*{Your turn}
\begin{itemize}
\raggedright
  \item \emph{Say it:} \emph{shonāno}
  \item \emph{Say it:} \emph{shonāno}, once more
  \item \emph{Say it:} \emph{didimā āmāder rājā ār rānir gôlpo shonālen.} — Grandmother told us the story of a king and a queen.
  \item \emph{Recall:} say the Bengali for to verify, then the Bengali for a tale, then say \emph{shonāno} again
\end{itemize}

\begin{tcolorbox}[breakable,colback=teal!4,colframe=teal!35!black,title={Before you move on}]
What does \bn{শোনানো} mean? (\textquotedblleft{}To tell a story to; to recite\textquotedblright{}.) Say it once more.
\end{tcolorbox}
\section[\texorpdfstring{\bn{বর্ণনা} \bn{করা} (bôrnônā kôrā)}{বর্ণনা করা (bôrnônā kôrā)}]{\bn{বর্ণনা} \bn{করা} (bôrnônā kôrā) — to describe}

\label{lesson:BN-C302-bornona-kora}



\begin{quote}
Before the new one: say the Bengali for a tale, then the Bengali for to tell a story.
\end{quote}

\subsection*{You'll want to know: \bn{বর্ণনা} \bn{করা}}
\textbf{\bn{বর্ণনা} \bn{করা}} — \emph{bôrnônā kôrā} — \textquotedblleft{}to describe\textquotedblright{}.

\begin{quote}
\bn{সাক্ষী} \bn{চোরের} \bn{চেহারা} \bn{বর্ণনা} \bn{করল।}
\end{quote}

\begin{quote}
\emph{sākkhī chorer chehārā bôrnônā kôrlô.}
\end{quote}

\begin{quote}
The witness described what the thief looked like.
\end{quote}

Past (he/she …): \bn{তিনি} \bn{জায়গাটা} \bn{বর্ণনা} \bn{করলেন}, \emph{tini jāygāṭā bôrnônā kôrlen}.

\begin{cousinweb}
\textbf{\bn{বর্ণনা} \bn{করা}} (\emph{bôrnônā kôrā}) puts Sanskrit \emph{varṇanā}, which Monier-Williams gives as 'praise, commendation', before \bn{করা}; it is built on \emph{varṇa}, 'colour', and to describe is to give a thing its colours.
\end{cousinweb}

\begin{grammarlens}[title={What you've built}]
One more word for telling what happened.
\end{grammarlens}

\subsection*{Your turn}
\begin{itemize}
\raggedright
  \item \emph{Say it:} \emph{bôrnônā kôrā}
  \item \emph{Say it:} \emph{bôrnônā kôrā}, once more
  \item \emph{Say it:} \emph{sākkhī chorer chehārā bôrnônā kôrlô.} — The witness described what the thief looked like.
  \item \emph{Recall:} say the Bengali for a tale, then the Bengali for to tell a story, then say \emph{bôrnônā kôrā} again
\end{itemize}

\begin{tcolorbox}[breakable,colback=teal!4,colframe=teal!35!black,title={Before you move on}]
What does \bn{বর্ণনা} \bn{করা} mean? (\textquotedblleft{}To describe\textquotedblright{}.) Say it once more.
\end{tcolorbox}
\section[\texorpdfstring{\bn{ঘটানো} (ghôṭāno)}{ঘটানো (ghôṭāno)}]{\bn{ঘটানো} (ghôṭāno) — to bring about; to cause}

\label{lesson:BN-C302-ghotano}



\begin{quote}
Before the new one: say the Bengali for to tell a story, then the Bengali for to describe.
\end{quote}

\subsection*{You'll want to know: \bn{ঘটানো}}
\textbf{\bn{ঘটানো}} — \emph{ghôṭāno} — \textquotedblleft{}to bring about; to cause\textquotedblright{}.

\begin{quote}
\bn{একটা} \bn{ছোট} \bn{ভুল} \bn{বড়} \bn{দুর্ঘটনা} \bn{ঘটাল।}
\end{quote}

\begin{quote}
\emph{ekṭā chhoṭo bhul bôṛo durghôṭônā ghôṭālô.}
\end{quote}

\begin{quote}
One small mistake caused a big accident.
\end{quote}

Past (he/she …): \bn{সেই} \bn{ভুল} \bn{দুর্ঘটনা} \bn{ঘটাল}, \emph{sei bhul durghôṭônā ghôṭālô}.

\begin{cousinweb}
\textbf{\bn{ঘটানো}} (\emph{ghôṭāno}) is \bn{ঘটা}, to happen — Sanskrit \emph{ghaṭatē}, 'is busy with; happens' (CDIAL 4407) — with the -ānō ending: to make something happen.
\end{cousinweb}

\begin{grammarlens}[title={What you've built}]
One more word for telling what happened.
\end{grammarlens}

\subsection*{Your turn}
\begin{itemize}
\raggedright
  \item \emph{Say it:} \emph{ghôṭāno}
  \item \emph{Say it:} \emph{ghôṭāno}, once more
  \item \emph{Say it:} \emph{ekṭā chhoṭo bhul bôṛo durghôṭônā ghôṭālô.} — One small mistake caused a big accident.
  \item \emph{Recall:} say the Bengali for to tell a story, then the Bengali for to describe, then say \emph{ghôṭāno} again
\end{itemize}

\begin{tcolorbox}[breakable,colback=teal!4,colframe=teal!35!black,title={Before you move on}]
What does \bn{ঘটানো} mean? (\textquotedblleft{}To bring about; to cause\textquotedblright{}.) Say it once more.
\end{tcolorbox}
\section[\texorpdfstring{\bn{ক্রমশ} (krômôsh)}{ক্রমশ (krômôsh)}]{\bn{ক্রমশ} (krômôsh) — gradually; little by little}

\label{lesson:BN-C302-kromosh}



\begin{quote}
Before the new one: say the Bengali for to describe, then the Bengali for to bring about.
\end{quote}

\subsection*{You'll want to know: \bn{ক্রমশ}}
\textbf{\bn{ক্রমশ}} — \emph{krômôsh} — \textquotedblleft{}gradually; little by little\textquotedblright{}.

\begin{quote}
\bn{ঝড়টা} \bn{ক্রমশ} \bn{দুর্বল} \bn{হল।}
\end{quote}

\begin{quote}
\emph{jhôṛṭā krômôsh durbôl hôlô.}
\end{quote}

\begin{quote}
The storm gradually grew weaker.
\end{quote}

\begin{cousinweb}
\textbf{\bn{ক্রমশ}} (\emph{krômôsh}) is Sanskrit \emph{kramaśas}, 'gradually, by degrees' (Monier-Williams), built on \emph{krama}, 'a step', the word in \textbf{\bn{ভাগ্যক্রমে}}.
\end{cousinweb}

\begin{grammarlens}[title={What you've built}]
One more word for telling what happened.
\end{grammarlens}

\subsection*{Your turn}
\begin{itemize}
\raggedright
  \item \emph{Say it:} \emph{krômôsh}
  \item \emph{Say it:} \emph{krômôsh}, once more
  \item \emph{Say it:} \emph{jhôṛṭā krômôsh durbôl hôlô.} — The storm gradually grew weaker.
  \item \emph{Recall:} say the Bengali for to describe, then the Bengali for to bring about, then say \emph{krômôsh} again
\end{itemize}

\begin{tcolorbox}[breakable,colback=teal!4,colframe=teal!35!black,title={Before you move on}]
What does \bn{ক্রমশ} mean? (\textquotedblleft{}Gradually; little by little\textquotedblright{}.) Say it once more.
\end{tcolorbox}
